% problem-000264 6 储氢技术与体积密度
\section*{6. 储氢技术与体积密度}
利用绿氢替代化石燃料实现工业、发电、交通等高排放领域深度脱碳是我国实现2030年碳达峰、2060年碳中和的“双碳目标”核心支撑路径之一。氢的高效存储是当前氢能产业面临的重大挑战。本题探讨几种储氢方法与体积储氢密度的关系。

6.1 高压储氢。压缩氢气是目前最常用的储氢方法之一。氢气被存储在压力维持在 350\~700 bar 的容器中，高压下氢气行为偏离理想气体，更适合用范德华方程 $(p+a/V_{\mathrm{m}}^{2})(V_{\mathrm{m}}-b)=RT$ 描述，其中 $V_{m}$ 为气体的摩尔体积，p 为气体压强。对氢气而言，方程中的 $a=0.02476\ J\ m^{3}\ mol^{-2}$ ， $b=2.661\times10^{-5}\ m^{3}\ mol^{-1}$ 。提示：此小题要求有效数字。

6.1.1 293 K 下，要达到 $20.0 \, kg \cdot m^{-3}$ 的储氢密度，所需的压力为多少 bar? 解答：

20.0 kg m $^{-3}$ 对应的 $V_{m}=2.016\times10^{-3}/20.0=1.008\times10^{-4}(m^{3}\ mol^{-1})$

(1 分)

$
p = \frac {R T}{V _ {m} - b} - \frac {a}{V _ {m} ^ {2}} = \frac {8 . 3 1 4 \times 2 9 3}{1 . 0 0 8 \times 1 0 ^ {- 4} - 2 . 6 6 1 \times 1 0 ^ {- 5}} - \frac {0 . 0 2 4 7 6}{(1 . 0 0 8 \times 1 0 ^ {- 4}) ^ {2}}
$

$
= 3. 0 4 \times 1 0 ^ {7} (\mathrm{Pa}) = 3. 0 4 \times 1 0 ^ {2} (\mathrm{bar})
$

(1 分)

给分原则：

$V_{m}$ 正确 1 分，答案正确 1 分，答案不用 bar 为单位不得分， $(3.03 \sim 3.07) \times 10^{2} \text{ bar 结果都算正确。结果保留三位有效数字，中间过程可保留四位有效数字。过程正确，答案有效数字错误，扣 1 分，无代入数据扣 1 分。}$

6.1.2 293 K 下，要想将储氢密度由 $20.0 \, kg m^{-3}$ 提高至 $40.0 \, kg m^{-3}$ ，所需压力（）:

(a) 为原来的 2 倍；(b) 低于原来的 2 倍；(c) 高于原来的 2 倍；(d) 无法确定。解答：

c (1 分),

\subsection*{给分原则：其他答案不得分}

6.1.3 工业上储氢罐的压强很少高于 700 bar，简要分析其原因。

解答:

储氢密度增加的倍数较压强增加的倍数小， (2 分)

700 Bar 后，增加压强对增加储氢密度的影响不大。从安全和成本的角度，工业上选择 350～700 Bar 的储气罐，不再增加压强。

\subsection*{给分原则：}

答出进一步提高压力对储氢密度帮助不大可得 2 分，如果只答出安全和成本/避免氢气泄漏等因素，只得 1 分。其他答案不得分

6.2 液相储氢。 $H_{2}$ 经液化后可在 1\~4 bar 的压力下稳定存储，但系统需维持极低温度。已知氢气的三相点为 $(7.041\ \text{kPa}, -259.35\ ^{\circ}\text{C})$ ，临界点为 $(12.86\ \text{bar}, -240.21\ ^{\circ}\text{C})$ 。

6.2.1 可能观察到液态氢的温度有: (a) 16 K; (b) 25 K; (c) 77 K; (d) 293 K。

解答:

a、b

给分原则：

每个 1 分；共 2 分。有错误选项就得 0 分。

6.2.2 计算 $\mathrm{H}_{2}$ 在 $27.15\mathrm{K}$ 液化所需的压力，并说明计算过程采取了哪些近似。解答：

由 Clausius-Clapeyron 方程： $\ln\frac{p_{2}}{p_{1}}=-\frac{\Delta vapH_{m}}{R}\left(\frac{1}{T_{2}}-\frac{1}{T_{1}}\right)$ (1 分)

代入三相点与临界点的数据，有

$
\ln {\frac {1 2 8 6}{7 . 0 4 1}} = - \frac {\Delta_ {v a p} H _ {m}}{8 . 3 1 4} \left(\frac {1}{2 7 3 . 1 5 - 2 4 0 . 2 1} - \frac {1}{2 7 3 . 1 5 - 2 5 9 . 3 5}\right)
$

解得 $\Delta_{vap}H_{m}=1.0283\times10^{3}\ (J\ mol^{-1})$ (2分)

以三相点数据计算，假设 $p_{3}$ 时，氢气的沸点为 27.15 K，则有

$
\ln \frac {p _ {3}}{7 . 0 4 1} = - \frac {1 . 0 2 8 \times 1 0 ^ {3}}{8 . 3 1 4} \left(\frac {1}{2 7 . 1 5} - \frac {1}{(2 7 3 . 1 5 - 2 5 9 . 3 5)}\right)
$

解得 $p_{3}=577.5\ (kPa)=5.775\ (bar)$ (1分)

或者用临界点数据计算， $\ln\frac{1286\ kPa}{p_{3}}=-\frac{\Delta_{vap}H_{m}}{8.314\ J\ K^{-1}\ mol^{-1}}\left(\frac{1}{(273.15-240.21)\ K}-\frac{1}{27.15\ K}\right)$

$
p _ {3} = 5 7 7. 4 \mathrm{kPa} = 5. 7 7 4 \mathrm{bar}
$

近似包括：将氢气视为理想气体（1 分）；认为蒸发热与温度无关（1 分）。

\subsection*{给分原则：}

● 本小问共 6 分，其中 Clausius-Clapeyron 方程书写 1 分，气化焓 2 分(代入数值 1 分，结果 1 分)，正确压力 1 分，近似 2 分。如果有跳步，结果正确得相应步骤对应得满分，错误则不得分。

● 最终答案 5.76 \~ 5.78 bar 范围内均给满分

\- 如果回答“近似认为Clausius-Clapeyron方程适用”，近似的2分只给1分。

6.2.3 若液氢的状态也可以近似用范德华方程描述，计算液氢密度的上限(单位： $\mathrm{kg}\ \mathrm{m}^{-3}$ )。
解答：

由范德华方程， $(p+a/V_{\mathrm{m}}^{2})(V_{\mathrm{m}}-b)=RT$

当 p 无穷大时， $V_{m}$ 最小值约等于范德华方程中的 b 值。(1 分)

因此，密度上限为 $\rho = \frac{M}{b} = \frac{2.016\times 10^{-3}\mathrm{kg mol^{-1}}}{2.661\times 10^{-5}\mathrm{m}^3\mathrm{mol}^{-1}} = 75.8\mathrm{kg m^{-3}}$ 。 (1分)

给分原则：本小问共2分，得出 $V_{\mathrm{m}}$ 最小值约等于 $b$ 值1分，结果1分，计算储氢密度时无数据代入扣1分。如没有说明 $V_{\mathrm{m}}$ 最小值约等于 $b$ 值，结果正确得满分，结果错误不得分。

6.3 有机液态储氢。有机液态储氢载体(LOHCs)具有高安全、易运输的优点。N-乙基咔唑(NEC)及其全氢化产物12H-NEC之间可以可逆转变，是极具应用前景的LOHC体系。

6.3.1 12H-NEC 的式量为 207.36，其密度 $\rho$ (单位：g cm $^{-3}$ ) 同温度 T (单位：K) 的关系为： $\rho = 1.1482329 - 0.0007092T$ ，计算 293.0 K 时 12H-NEC 的储氢密度为多少 kg m $^{-3}$ ?

（图：）

（图：）
NEC
12H-NEC

解答:

293 K 时，12H-NEC 的密度为

$
\rho / (\mathrm{gcm} ^ {- 3}) = 1. 1 4 8 2 - 0. 0 0 0 7 0 9 2 \times 2 9 3. 0 = 0. 9 4 0 4
$

12H-NEC 的分子量为 207.357，则

等效贮氢密度=940.4 × $\frac{12\times1.008}{207.36}=54.86\ (kg\ m^{-3})$

给分原则：公式和代入数值1分，答案准确1分。

（图：）

6.3.2 12H-NEC 脱氢时，会生成 8H-NEC (产物 1) 和 4H-NEC (产物 2) 两种中间产物，分别是 NEC 加上 8 个和 4 个氢原子的化合物。分别画出产物 1 和产物 2 的结构，不要求立体化学。
解答：

（图：）

产物1

(1 分); 产物 2

给分原则：其他答案不得分

(1 分)

6.3.3 12H-NEC 存在多少种立体异构体(包含对映异构)?

解答:

10 种 (2 分)

6.3.4 已知 12H-NEC 脱氢到 NEC 是较强的吸热反应，画出脱氢温度最低的异构体结构。解答：

（图：）

(2 分)

给分原则：

结构需体现出五元环上四个氢为全顺式，其他答案不得分

6.4 多孔吸附储氢材料。MOF-5 的比表面积高达数千平方米每克，可以通过表面吸附储氢。若氢气在 MOF-5 表面的吸附过程， $\Delta H_{ads} = -5.20 \, kJ mol^{-1}$ ， $\Delta S_{ads} = -80.3 \, J mol^{-1} K^{-1}$ ，均视为同温度无关。表面吸附量 $\Gamma$ 与 $H_{2}$ 压力 p 服从 Langmuir 吸附等温式： $\Gamma/\Gamma_{\infty} = Kp/(1 + Kp)$ ，其中 K 为表面吸附反应的平衡常数， $\Gamma_{\infty}$ 为常数。

6.4.1 MOF-5 的密度为 $0.30 \, g cm^{-3}$ (包含结构中的孔)，77 K、100 bar 下 MOF-5 可以吸附 7.5 wt% 的氢气，计算 77 K 下 MOF-5 通过表面吸附能够达到的最大储氢密度 $(\mathrm{kg}\cdot\mathrm{m}^{-3})$ 。

解答:

MOF-5 (s) + H₂(g) = MOF-5 \~ H₂ (s)

77 K 时，

$
K _ {p} = \exp \left(- \frac {\varDelta_ {r} G _ {m} ^ {\emptyset}}{R T}\right) = \exp \left(- \frac {\varDelta_ {r} H _ {m} ^ {\emptyset} - T \varDelta_ {r} S _ {m} ^ {\emptyset}}{R T}\right) = \exp \left(- \frac {- 5 2 0 0 - 7 7 \times (- 8 0 . 3)}{8 . 3 1 4 \times 7 7}\right) = 0. 2 1 5\tag{2 分}
$

$
\frac {\Gamma}{\Gamma_ {\infty}} = \frac {K _ {p} p}{1 + K _ {p} p}, \text {即有} \frac {0 . 0 7 5}{\Gamma_ {\infty}} = \frac {0 . 2 1 5 p / p ^ {\emptyset}}{1 + 0 . 2 1 5 p / p ^ {\emptyset}} = \frac {0 . 2 1 5 \times 1 0 0}{1 + 0 . 2 1 5 \times 1 0 0}, \Gamma_ {\infty} = 7.85\tag{2 分}
$

储氢密度为 $0.30 \times 10^{3} \mathrm{~kg} \mathrm{~m}^{-3} \times 7.85\% = 24 (\mathrm{kg} \mathrm{~m}^{-3})$

(1 分)

给分原则：

本小问共 5 分， $K_{p}$ 计算过程中，代入数值 1 分，答案正确 1 分，共 2 分， $\Gamma_{\infty}$ 计算过程中，代入数值 1 分，答案正确 1 分，共 2 分，储氢密度数值正确 1 分，一处无数据代入扣 1 分。此处 $7.5 \mathrm{wt}\%$ 是吸附氢气和 MOF-5 质量之比，如果计算成 $(\mathrm{MOF} - 5 + \mathrm{H}_{2})$ 的 $7.5 \mathrm{wt}\%$ ，按照计算错误处理。如果有跳步，结果正确得相应步骤对应的分数，错误则不得分。

6.4.2 除表面吸附外，MOF-5 的孔中还可以存储 $\mathrm{H}_{2}$ ，MOF-5 的比孔容为 $1.27\mathrm{cm}^{3}\mathrm{g}^{-1}$ 。向一个 $1.00\mathrm{L}$ 的储氢罐中装填满 MOF-5 材料后，计算 $77\mathrm{K}$ 、 $100\mathrm{bar}$ 下其储氢量为未装填 MOF-5 材料的多少倍？假设 MOF-5 表面吸附 $\mathrm{H}_{2}$ 后其孔体积不变，孔中的 $\mathrm{H}_{2}$ 符合范德华方程。在 $77\mathrm{K}$ ， $100\mathrm{bar}$ 下的 $V_{\mathrm{m}} = 6.853 \times 10^{-5}\mathrm{m}^{3}\mathrm{mol}^{-1}$ 。

\subsection*{解答:}

1 L MOF-5 孔的体积为:

$
V = m _ {\mathrm{MOF}} \cdot \text { pore } \% = \rho_ {\mathrm{MOF}} V \cdot \text { pore } \% = 0. 3 0 \times 1 0 0 0 \times 1. 2 7 = 3 8 1 (\mathrm{cm} ^ {3})\tag{1 分}
$

孔中能容纳的氢气的质量为:

$
m _ {\text { pore }} = M _ {r} \left(\mathrm{H} _ {2}\right) \cdot V / V _ {\mathrm{m}} = 2. 0 1 6 \times 3 8 1 \times 1 0 ^ {- 6} / (6. 8 5 3 \times 1 0 ^ {- 5}) = 1 1. 2 1 (\mathrm{g})\tag{1 分}
$

吸附储氢量:

$
m _ {\mathrm{abs}} = m _ {\mathrm{MOF}} \cdot \mathrm{H} - \mathrm{wt} \% = \rho_ {\mathrm{MOF}} V \cdot \mathrm{H} - \mathrm{wt} \% = 0. 3 0 \times 1 0 0 0 \times 7. 5 \% = 2 2. 5 (\mathrm{g})
$

总储氢量 $m_{total} = m_{pore} + m_{abs} = 11.21 + 22.5 = 33.7 \, (g)$

(1 分)

得出 MOF “吸附+孔”的储氢总量，1 分

1 L 空罐体的储氢量

$
m _ {\mathrm{H}} = M _ {r} \left(\mathrm{H} _ {2}\right) \cdot n \left(\mathrm{H} _ {2}\right) = M _ {r} \left(\mathrm{H} _ {2}\right) \cdot V / V _ {m} = 2. 0 1 6 \times 1. 0 0 \times 1 0 ^ {- 3} / (6. 8 5 3 \times 1 0 ^ {- 5}) = 2 9. 4 (\mathrm{g})\tag{1 分}
$

得出空罐的储氢总量，1 分

因此加入 MOF 导致的储氢量为空罐的 33.7/29.4 = 1.15 倍

(1 分)

\subsection*{给分原则:}

本小问共 5 分，得出 MOF 的储氢总量 3 分（其中孔体积 1 分，孔中的 $H_{2}$ 量 1 分，MOF-5 储氢总量 1 分），空罐储氢质量 1 分，比例 1 分。一处无数据代入扣 1 分。如果有跳步，结果正确得到相应步骤对应的分数，错误则不得分。
